% concatenated from toArxivMay11.tex
 
\documentclass{article} 

\usepackage{amsmath}  
\usepackage{amsthm}


\theoremstyle{plain}
\newtheorem{thm}{Theorem}
\newtheorem{lem}[thm]{Lemma}
\newtheorem{cor}[thm]{Corollary}
\newtheorem{prop}[thm]{Proposition}

\theoremstyle{definition}
\newtheorem{defn}{Definition}
\newtheorem{exa}{Example}

\title{Chromatic numbers from edge ideals:\\ Graph classes with vanishing syzygies are polynomially $\chi$-bounded}
\author{ Alexander Engström\\ alexander.engstrom@fra.se }
\date\today

\begin{document}
\maketitle

\begin{abstract}
The chromatic number $\chi$ of a graph is bounded from below by its clique number $\omega,$ but it can be arbitrary large. Perfect graphs are defined by $\chi=\omega$ for all induced subgraphs. An interesting relaxation are $\chi$-bounded graph classes, where $\chi\leq f(\omega).$ It is not always possible to achieve this with a polynomial $f.$


The edge ideal $I_G$ of a graph $G$ is generated by monomials $x_ux_v$ for each edge $uv$ of $G.$ The bi-graded betti numbers $\beta_{i,j}(I)$ are central algebraic geometric invariants. We study the graph classes where for some fixed $i,j$ that syzygy vanishes, that is, $\beta_{i,j}(I_G)=0.$ 

We prove that $\chi\leq f(\omega),$ where $f$ is a polynomial of degree $2j-2i-4.$
For the elementary special case $\beta_{i,2i+2}(I_G)=0,$ this  amounts to that $(i+1)K_2$--free graphs are ${\omega-1+2i \choose 2i}$--colorable, improving on an old combinatorial result by Wagon. 
We also show that triangle-free graphs with $\beta_{i,j}(I_G)=0$ are $(j-1)$--colorable.

Complexity wise, we show that these colorings can be derived in time $O(n^3)$ for graphs on $n$ vertices. Moreover, we show that for almost all graphs with parabolic $i,j,$ there are better bounds on $\chi.$
\end{abstract}

\section{Introduction}

For a graph $G,$ the largest number of vertices in a clique is a natural lower bound for its chromatic number: $\omega(G) \leq \chi(G).$ In general that inequality may be arbitrary bad, as there are triangle-free graphs of any chromatic number. On the other hand, a perfect graph is characterized by that $\omega(H) = \chi(H)$ for all of its induced subgraphs $H.$ The last decades have seen an intense activity in extending our understanding of how to measure and categorize the perfectness of graphs. A class of graphs is \emph{$\chi$--bounded} by the function $f$ if $\chi(G) \leq f(\omega(G))$ for all graphs $G$ in the class. The central Gyárfás-Sumner conjecture, which has been open for fifty years, is that for any forest $F,$ the $F$--free graphs are $\chi$--bounded.
For established $\chi$--bounded graph classes the common next question is if the bounding function is polynomial or not. In general it can be arbitrarily bad \cite{REF1}. Examples of polynomially $\chi$--bounded graph classes are even-hole-free graphs \cite{CS23}, circle graphs \cite{DM21, Dav22} and bounded twin-width graphs \cite{BT25}.

For a graph $G$ on vertex set $V$ and edge set $E,$ its edge ideal in $\mathbf{k}[x_v \mid v\in V]$ is generated by the monomials $x_ux_v$ for $uv \in E(G).$
In discussing betti numbers $\beta_{i,j}(I_G)$ of edge ideals, we assume from now that $j\geq 2i+2,$ because otherwise they are trivially zero. For general combinatorial considerations regarding edge ideals, see for example \cite{de} or \cite{EO23}.
From a geometric stand point, the first interesting case is if almost all $\beta_{i,j}(I_G)$ vanishes, except for if $j=i+2.$ Then the ideal is said to have a \emph{linear} resolution. In a seminal paper building a bridge between algebra and combinatorics, Fr\"oberg proved \cite{fro} that an edge ideal has a linear resolution if and only if the graph is the complement of a chordal graph. The complement of a chordal graph is perfect, setting the backdrop for this paper. In Fröberg's setting an infinite number syzygyies were set to vanish, but we study when only a single one does. The last years a number of interesting combinatorial results have been established with this methodology, for example Randriambololona \cite{ran} won the best paper award of Eurocrypt 2025 investigating the post-quantum cryptographic algorithm McEliece.

The main result of this paper is proved in Section 4: Corollary \ref{mainCor1} states that the class of graphs with $ \beta_{i,j}(I_G) = 0$ is $\chi$--bounded by
\[ {\omega-1+2(j-i-2) \choose 2(j-i-2)} +  {j-2 \choose 2(j-i-2)+1}, \] 
which is a polynomial in $\omega$ of degree $2(j-i-2).$ For triangle-free graphs the slightly stronger result  $\chi(G) \leq j-1$ is proved in Corollary \ref{mainCor2}.

In Section 5 we show in Corollary \ref{asymBound} that if $(j-i)^2 \geq j+i+2$, then almost all graphs $G$ with $ \beta_{i,j}(I_G) = 0$ are $\chi$-bounded by
\[ {\omega-1+2(j-i-2) \choose 2(j-i-2)}. \] 

In the penultimate Section 6 we show in Corollary \ref{n3} that the graph colorings in our main result may be done in time $O(n^3)$ for graphs on $n$ vertices. For the triangle-free case, it is in time $O(n^2)$ according to Corollary \ref{n2}. Some further considerations on shorter proofs with worse chromatic number bounds are recorded in the final Section 7.

Before reaching the main part of our paper, some elementary considerations on $\chi$-boundedness are recorded in Section 2. It is  included due to the way more technical versions of the same arguments in Section 3 that aims towards applications in the algebraic setting. Although it is all fairly elementary, it includes an improvement of an old influential result by Wagon \cite{W80}, proving that $(i+1)K_2$--free graphs are ${\omega-1+2i \choose 2i}$--colorable.

\section{Ordinary polynomially $\chi$--bounded graph classes}

The chromatic number of a graph $G$ is denoted by $\chi(G).$ The order of the largest clique of $G$ is denoted by $\omega(G).$

\begin{defn} An hereditary class of graphs is \emph{$\chi$-bounded by $f$} if $\chi(G)\leq f(\omega(G))$ for all graphs $G$ in the class.
\end{defn}

The disjoint union of two graphs $G$ and $H$ is denoted $G \cup H$ and multiples by $pG=\cup_p G.$
The join $G+H$ is $G \cup H$ with edges added between all of $G$ and $H.$ For a graph $G$ and a set $\bf{H}$ of graphs, let $ G \cup {\bf H  } = \{ G \cup H \mid H \in {\bf H} \}$ and  $ G + {\bf H  } = \{ G + H \mid H \in {\bf H} \}.$ The induced subgraphs on $U\subseteq V(G)$ of $G$ is denoted $G[U].$ Basic graphs on $n$ vertices are the complete graph $K_n,$ the cycle $C_n$ and the path $P_n.$

Wagon \cite{W80} proved the following influential result:
\begin{prop}\label{wagon} If $f_1^{\mathrm{W}}(\omega)=1$ and $f_{n+1}^{\mathrm{W}}(\omega)={ \omega \choose 2} f_n^{\mathrm{W}}(\omega) + \omega,$ then the class of $pK_2$--free graphs is $\chi$--bounded by $f_p^{\mathrm{W}}(\omega).$ In particular, the class of $2K_2$--free graphs is $\chi$--bounded by ${\omega +1 \choose 2}.$
\end{prop}
The proof of Lemma \ref{K2extend} below builds on the proof of Proposition \ref{wagon} by Wagon, who attributes the basic idea to Erdös and Hajnal. The main difference between the proofs is that the chromatic number for some induced subgraphs $G[C_{i,j}]$ are better estimated using an idea by Bharati and Chodum \cite{BC18}. But the main point of Lemma \ref{K2extend} is actually that it serves as a model for the more technical proof of Theorem \ref{bigthm}, which is crucial to derive the main results of this paper.

\begin{lem}\label{K2extend}
If ${\bf H}$ is a set of graphs and the class of ${\bf H}$--free graphs is  $\chi$-bounded by $f_{\bf H},$ then the class of  $K_2 \cup {\bf H  } $--free graphs is $\chi$-bounded by
\[
f_{K_2 \cup {\bf H  } }(\omega) = \sum_{k=1}^{\omega}  (\omega-k+1)  f_{{\bf H  } }(k).
\]
\end{lem}

\begin{proof}
Let $G$ be a $K_2 \cup {\bf H  } $--free graph with a clique on the vertices $1,2,\ldots, \omega=\omega(G).$
For $1\leq i < j \leq \omega$ define
\[
C_{i,j}' = \{ v\in V(G) \setminus \{1,2,\ldots, \omega\} \mid vi,vj \not\in E(G).  \}
\]
For any induced subgraph $H$ of $C_{i,j}'$, the graph $K_2 \cup H$ on $\{i,j\} \cup V(H)$ is an induced subgraph of $G.$ By assumption, $G$ is a $K_2 \cup {\bf H  } $--free graph, and, thus, the induced subgraph of $G$ on $C_{i,j}' $ is an  ${\bf H  } $--free graph.

Set
\[
C_{i,j} = C_{i,j}' \setminus  \left(\,\, \bigcup_{k=1}^{i-1} \bigcup_{l=k+1}^\omega C_{k,l}'   \quad \cup \quad  \bigcup_{l=i+1}^{j-1} C_{i,l}'  \,\,  \right). 
\]
If $v\in C_{i,j}$ and $1\leq k < i$ then $kv \in E(G)$ as 
\[ v\in C_{i,j}' \Rightarrow iv \not\in E(G) \textrm{ and }  jv \not\in E(G) \]
and
\[ v\not\in C_{k,j}' \Rightarrow kv \in E(G) \textrm{ or }  jv \in E(G). \]
If $v\in C_{i,j}$ and $i< l < j$ then $lv \in E(G)$ as 
\[ v\in C_{i,j}' \Rightarrow iv \not\in E(G) \textrm{ and }  jv \not\in E(G) \]
and
\[ v\not\in C_{i,l}' \Rightarrow iv \in E(G) \textrm{ or }  lv \in E(G). \]
It follows that any clique in $G[C_{i,j}]$ extends to a larger clique in $G$ by the $j-2$ vertices
$\{1,2,\ldots,i-1\} \cup \{ i+1,i+2,\ldots, j-1 \}.$ 
Thus
\[ \omega(G[C_{i,j}]) \leq \omega(G)-j+2\]
and
\[ \chi(G[C_{i,j}]) \leq f_{{\bf H  } }( \omega(G)-j+2 ). \]

So far we have partitioned the vertices in $G$ outside the clique on $1,2,\ldots, \omega=\omega(G)$  that has at least two non-edges to the clique. There cannot be a vertex with no non-edges to the clique, as the clique is maximal in $G.$ 
For $1\leq i \leq \omega(G)$ set
\[ D_i'
= \left\{ v\in V(G) \setminus \{1,2,\ldots, \omega\} \mid vi \not\in E(G); \{v1,v2, \ldots, v\omega \} \setminus vi \subseteq E(G) 
%v1,v2,\ldots v(i-1),v(i+1),\ldots v\omega 
 \right \}
\]
If there would be an edge between two vertices $u$ and $v$ in $D_i',$ then the vertices
\[ 1,2, \ldots, i-1, i+1, \ldots, \omega, u, v \]
would constitute a clique on $\omega+1$ vertices, which is a contradiction. Thus, $D_i'$ is independent. As there are no edges from $i$ to vertices of $D_i',$ the set
\[ D_i = \{i\} \cup D_i' \]
is independent for all $i.$ We only need one color for an independent set, and
\[ \chi(G[D_i]) = 1 \leq  f_{{\bf H  } }( 1). \]
Coloring each part of the partition
\[ V(G) \quad = \quad   \bigcup_{i=1}^\omega D_{i}      \quad \cup \quad    \bigcup_{1\leq i < j \leq \omega}  C_{i,j}  \]
gives
\[\begin{array}{rcl} 
\chi(G) & \leq & \displaystyle  \sum_{i=1} ^\omega \chi(G[D_i]) +  \sum_{1\leq i < j \leq \omega}  \chi(G[C_{i,j}]) \\
& \leq & \displaystyle  \sum_{i=1} ^\omega  f_{{\bf H  } }( 1) +  \sum_{1\leq i < j \leq \omega}   f_{{\bf H  } }( \omega-j+2 ) \\
& = & \displaystyle  \omega  f_{{\bf H  } }( 1) +  \sum_{j=2}^\omega   (j-1)f_{{\bf H  } }( \omega-j+2 ) \\
& = & \displaystyle  \sum_{j=2}^{\omega+1}   (j-1)f_{{\bf H  } }( \omega-j+2 ) \\
& = & \displaystyle \sum_{k=1}^{\omega}  (\omega-k+1)  f_{{\bf H  } }(k)
\end{array}
\]
\end{proof}

This is completely elementary 
\begin{lem}\label{doubleBumpLemma}
For non-negative integers $n$ and $m,$
\[ 
\sum_{i=0}^n   (n+1-i)    {i + m   \choose m } =  {n + m+2   \choose m+2}.
\]
\end{lem}
\begin{proof} The right hand side counts the number of order $m+2$ subsets of $\{0,1,\ldots,n+m+1\}.$ Put those subsets in buckets depending on the value of the second largest element of each subset. The value of the second largest element is $m+i$ for $0\leq i \leq n.$ There are $ {i + m   \choose m } $ ways to select the $m$ smallest elements and $ (n+1-i) $ ways to select the largest element. Summing up the sizes of the buckets gives the equality.
\end{proof}

\begin{thm}\label{pK2extend}
If ${\bf H}$ is a set of graphs and the class of ${\bf H}$--free graphs is  $\chi$-bounded by $f_{\bf H},$ then the class of $pK_2 \cup {\bf H  } $--free graphs is $\chi$-bounded by
\[
f_{pK_2 \cup {\bf H  } }(\omega) = \sum_{k=1}^{\omega}   {\omega-k + 2p-1  \choose 2p-1}    f_{{\bf H  } }(k).
\]
\end{thm}

\begin{proof}
For $p=1$ this is exactly Lemma \ref{K2extend}. Now we proceed by induction on $p>1.$ It is known that 
\[
f_{(p-1)K_2 \cup {\bf H  } }(\omega) = \sum_{k=1}^{\omega}   {\omega-k + 2(p-1)-1  \choose 2(p-1)-1}    f_{{\bf H  } }(k) =  \sum_{k=1}^{\omega}   {\omega-k + 2p-3  \choose 2p-3}    f_{{\bf H  } }(k) 
\]
is $\chi$-bounding the class of $(p-1)K_2 \cup {\bf H  }$--free graphs.
According to Lemma \ref{K2extend},
\[
\begin{array}{rcl}
f_{pK_2 \cup {\bf H  } }(\omega) &= & \displaystyle \sum_{k=1}^{\omega}  (\omega-k+1)  f_{ (p-1)K_2 \cup  {\bf H  } }(k) \\
&= & \displaystyle \sum_{k=1}^{\omega}  (\omega-k+1)   \sum_{l=1}^{k}   {k-l + 2p-3  \choose 2p-3}    f_{{\bf H  } }(l) \\
&= & \displaystyle \sum_{k=1}^{\omega} \sum_{l=1}^{k}   (\omega-k+1)    {k-l + 2p-3  \choose 2p-3}    f_{{\bf H  } }(l) \\
&= & \displaystyle \sum_{l=1}^{\omega} \sum_{k=l}^{\omega}   (\omega-k+1)    {k-l + 2p-3  \choose 2p-3}    f_{{\bf H  } }(l) \\
\end{array}
\]
is $\chi$-bounding the class of $pK_2 \cup {\bf H  }$--free graphs. According to Lemma \ref{doubleBumpLemma} with  $m=2p-3, $ and $i=k-l$ and $n=\omega-l,$
\[ 
\sum_{k=l}^{\omega}   (\omega-k+1)    {k-l + 2p-3  \choose 2p-3} =  {\omega-l + 2p-1  \choose 2p-1}
\]
and
\[
\begin{array}{rcl}
f_{pK_2 \cup {\bf H  } }(\omega) &= &  \displaystyle \sum_{l=1}^{\omega}  {\omega-l + 2p-1  \choose 2p-1}    f_{{\bf H  } }(l) \\
\end{array}
\]
\end{proof}


Another elementary lemma:

\begin{lem}\label{singleBumpLemma}
$ \sum_{i=0}^n {m+i \choose m} = { n+m+1 \choose m+1 }. $
\end{lem}
\begin{proof}
Consider selection of $m+1$ elements of $\{0,1,\ldots, n+m\}$ and bucket depending on the largest element and sum up each bucket. 
\end{proof}

This will be useful in the algebraic setting:

\begin{cor}\label{corConstF}
 If $f_{{\bf H  } }(1) = 1$ and   $f_{{\bf H  } }(\omega) = c_{\bf H  }$ for $\omega > 1$ is a $\chi$--bounding function for ${\bf H  }$--free graphs, then
\[ f_{pK_2 \cup {\bf H  } }(\omega) =  {\omega + 2p-2  \choose 2p-1}  +   c_{\bf H  }  {\omega + 2p-2  \choose 2p} \]
is a $\chi$--bounding function for ${ pK_2 \cup \bf H  }$--free graphs.
\end{cor}
\begin{proof}
According to Theorem \ref{pK2extend} and Lemma \ref{singleBumpLemma},
\[ \begin{array}{rcl}
f_{pK_2 \cup {\bf H  } }(\omega) & = & \displaystyle  {\omega-1 + 2p-1  \choose 2p-1}  +   c_{\bf H  }   \sum_{k=2}^{\omega}   {\omega-k+ 2p-1  \choose 2p-1}  \\
& = & \displaystyle  {\omega+ 2p-2  \choose 2p-1}  +   c_{\bf H  }   \sum_{i=0}^{\omega-2}   {i + 2p-1  \choose 2p-1}  \\
& = & \displaystyle  {\omega + 2p-2  \choose 2p-1}  +   c_{\bf H  }  {\omega + 2p-2  \choose 2p}.     \\
\end{array}
\]
\end{proof}

\begin{cor}\label{corBoundingSimple}
If ${\bf H}$ is the set of graphs on $n$ vertices of maximal degree $n-1,$ then 
\[ f_{pK_2 \cup {\bf H  } }(\omega) =  {\omega + 2p-2  \choose 2p-1}  +   (n-1)  {\omega + 2p-2  \choose 2p} \]
is a $\chi$--bounding function for ${ pK_2 \cup \bf H  }$--free graphs.
\end{cor}
\begin{proof}
The ${\bf H}$-free graphs are of maximal degree at most $n-2,$ and $(n-1)$--colorable. Use Corollary \ref{corConstF}.
\end{proof}

\begin{cor}\label{corInducedMatching} The function
\[
f_{pK_2}(\omega) =  {\omega-1 + 2(p-1)  \choose 2(p-1)}
\]
is $\chi$--bounding the $pK_2$--free graphs.
\end{cor}
\begin{proof} Use ${\bf H}=\{K_2\}$ and $ c_{\bf H  }=1$ in Corollary \ref{corConstF} with $(p-1)K_2$ to get
\[ f_{pK_2}(\omega) = {\omega + 2(p-1)-2  \choose 2(p-1)-1}  +   1  {\omega + 2(p-1)-2  \choose 2(p-1)} = {\omega-1 + 2(p-1)  \choose 2(p-1)}
\]
\end{proof}

Let's compare Corollary \ref{corInducedMatching} to Wagon's old result in Proposition \ref{wagon}. It is easy to prove that:
\begin{prop}
\[ \lim_{\omega \rightarrow \infty} \,\,   \frac{f_p^{\mathrm{W}}(\omega)}{\omega^{2(p-1)}} = \frac{1}{2^{p-1}} 
\quad
\textrm{and}
\quad
\lim_{\omega \rightarrow \infty} \,\,   \frac{f_{pK_2}(\omega)}{\omega^{2(p-1)}} = \frac{1}{(2p-2)!} 
\]
\end{prop}
So, for large $\omega,$ the new bound improves on Wagon's by a factor of $\frac{(2p-2)!}{2^{p-1}}.$
\begin{prop}
If $\omega>2$ and $p>2$ then
\[ f_p^{\mathrm{W}}(\omega) > f_{pK_2}(\omega) \]
and otherwise they are equal.
\end{prop}
\begin{proof} The recursion to prove Wagon's $f_p^{\mathrm{W}}(\omega)$ is a simpler version of the proof of Lemma \ref{K2extend}. Instead of the inequality
\[ \chi(G[C_{i,j}]) \leq f_{p-1}^{\mathrm{W}} ( \omega(G)-j+2 ) \]
the employed 
\[ \chi(G[C_{i,j}]) \leq f_{p-1}^{\mathrm{W}} ( \omega(G) ) \]
is worse when $\omega>2$ and $p>2$. That they are equal otherwise is straightforward.
\end{proof}

Here is another corollary of Theorem~\ref{pK2extend}.
\begin{cor} If $\bf H$ is a set of graphs such that all $\bf H$-free graphs are perfect, then
\[ f_{pK_2 \cup {\bf H  } }(\omega) = { \omega + 2p \choose 2p+1 } \]
is a $\chi$--bounding function for ${ pK_2 \cup \bf H  }$--free graphs.
\end{cor}
\begin{proof}
The $\bf H$-free graphs are $\chi$-bounded by $ f_{{\bf H  } }(k) = k$ as they are perfect. By Theorem~\ref{pK2extend} and Lemma~\ref{doubleBumpLemma},
\[
\begin{array}{rcl}
\displaystyle
f_{pK_2 \cup {\bf H  } }(\omega) & =& \displaystyle  \sum_{k=1}^{\omega}   {\omega-k + 2p-1  \choose 2p-1}    f_{{\bf H  } }(k)  \\ 
& = &  \displaystyle  \sum_{k=1}^{\omega}   {\omega-k + 2p-1  \choose 2p-1}  k \\
& = & \displaystyle  \sum_{i=0}^{\omega-1}  ((\omega-1)+1-i)   {i + 2p-1  \choose 2p-1}  \\
  & = &  \displaystyle { \omega - 1 + 2p -1 + 2\choose 2p-1 +2 } \\
  & = &  \displaystyle  { \omega + 2p \choose 2p+1 }
\end{array}
\]
is a $\chi$--bounding function for ${ pK_2 \cup \bf H  }$--free graphs.
\end{proof}

The following two results were proved by  Bharathi and Choudum \cite{BC18} in the case $p=1.$
\begin{cor}
The ${ pK_2 \cup P_4}$--free graphs are $\chi$-bounded by $\omega + 2p \choose 2p+1$.
\end{cor}
\begin{proof}
$P_4$--free graphs are perfect \cite{S74}.
\end{proof}
\begin{cor}
The ${ pK_2 \cup P_3}$--free graphs are $\chi$-bounded by $\omega + 2p \choose 2p+1$.
\end{cor}
\begin{proof}
$P_3$--free graph are disjoint unions of cliques, and thus perfect.
\end{proof}

For graphs that are almost perfect, similar results are possible:

\begin{cor}
The $\{ (p+2)K_2, pK_2 \cup diamond\}$--free graphs are $\chi$-bounded by $ {\omega + 2p \choose 2p+1} +   {\omega + 2p-3  \choose 2p-1}  $.
\end{cor}
\begin{proof}
Let ${\bf H} = \{2K_2, diamond\}.$
According to Theorem 4 of \cite{BC18}, the $\bf H$-free graphs are $\chi$-bounded by $ f_{{\bf H  } }(k) = k$ for all $k,$ with the exception $ f_{{\bf H  } }(2) = 3.$
By Theorem~\ref{pK2extend} and Lemma~\ref{doubleBumpLemma}, 
\[
\begin{array}{rcl}
\displaystyle
f_{pK_2 \cup {\bf H  } }(\omega) & =& \displaystyle  \sum_{k=1}^{\omega}   {\omega-k + 2p-1  \choose 2p-1}    f_{{\bf H  } }(k)  \\ 
& = &  \displaystyle   {\omega-2 + 2p-1  \choose 2p-1}  + \sum_{k=1}^{\omega}   {\omega-k + 2p-1  \choose 2p-1}  k \\
& = &  \displaystyle   {\omega-2 + 2p-1  \choose 2p-1}  + { \omega + 2p \choose 2p+1 }.
\end{array}
\]
\end{proof}

One step up from the perfect graphs are the perfectly divisible graphs. A graph $G$ is said to be \emph{perfectly divisible} if for all induced subgraphs $H$ of it, $V (H)$ can be partitioned into two sets $A$ and $B$ such that $H[A]$ is perfect and $\omega(H[B])<\omega(H).$ This lemma was proved by Chudnovsky and Sivaraman \cite{CS19}.
\begin{lem}
If $G$ is a perfectly divisible graph, then $\chi(G) \leq {\omega(G)+1 \choose 2}$.
\end{lem}
Along the lines above, it is straightforward to prove this this:
\begin{cor} If ${\bf H}$ is a set of graphs such that all ${\bf H}$--free graphs are perfectly divisible, then the $(p-1)K_2\cup {\bf H}$--free graphs are $\chi$--bounded by ${\omega-1 + 2p  \choose 2p}.$
\end{cor}
For example, Chudnovsky and Sivaraman \cite{CS19} proved that $\{P_5,bull\}$--free graphs are perfectly divisible.
\begin{cor} The $(p-1)K_2\cup  \{P_5,bull\}$--free graphs are $\chi$--bounded by ${\omega-1 + 2p  \choose 2p}.$
\end{cor}

\section{A class of graphs for vanishing syzygies}

Now we turn to the main consideration of this paper, how to deal with the induced subgraphs that are prohibited due to that syzygies vanish. 

\begin{defn}
Let ${\bf B}_{n,0}$ be the set of graphs on $n\geq 2$ vertices of maximal degree $n-1$, and let 
\[{\bf B}_{n,d} = \bigcup_{m=2d}^{n-2} K_{n-m-2}+(K_2\cup {\bf B}_{m,d-1}) \] for any integers $d>0$ and $n\geq 2(d+1).$
\end{defn}

\begin{exa}
The bowtie is in $B_{5,1}.$
\end{exa}

\begin{defn}
For any integers $d\geq 0$ and $n\geq 2(d+1)$ let the function $g_{n,d}: \{1,2,\ldots\} \rightarrow \{1,2,\ldots\}$ be
\[\begin{array}{rcll}
g_{n,d}(1) &= &1 & \textrm{for $d\geq 0$ and $n\geq 2(d+1),$} \\
g_{n,0}(\omega) &= & n-1 & \textrm{for $\omega > 1$ and $n\geq 2,$}\\
g_{n,d}(\omega) & = & \displaystyle  \sum_{k=1}^{\omega} (\omega-k+1) g_{\max(n+k-\omega-2,2d),d-1}( k ) &  \textrm{for $\omega > 1, d>0,$ and $n\geq 2(d+1).$}\\
\end{array}
\]
\end{defn}

\begin{thm}\label{bigthm}
For any integers $d\geq 0$ and $n\geq 2(d+1)$ the class of ${\bf B}_{n,d}$--free graphs is $\chi$--bounded by $g_{n,d}(\omega).$
\end{thm}

\begin{proof}
We are to show that $\chi(G) \leq g_{n,d}(\omega(G))$ for every ${\bf B}_{n,d}$--free graph $G.$

First note that for graphs with $\omega(G)=1$ it is satisfied, as $g_{n,d}(1)=1$ for all $n,d.$ In calculations later on in the proof, $g_{n,d}(1)$ will sometimes be replaced by $g_{\ast,\ast}(1),$ as the value of the function at $1$ is independent of the index $n,d.$

The proof will be by induction on $d.$ The base case is $d=0.$ By definition, ${\bf B}_{n,0}$ is the set of graphs on $n$ vertices of maximal degree $n-1.$ A ${\bf B}_{n,0}$--free graph $G$ contains no induced subgraph on 
$n$ vertices of maximal degree $n-1,$ or equivalently, the maximal degree of $G$ is at most $n-2.$ As the maximal degree of $G$ is at most $n-2,$ one may color it with $n-1$ colors, and $\chi(G) \leq n-1 = g_{n,0}(\chi(G))$ as desired.

Now to the induction step. We are provided with a ${\bf B}_{n,d}$--free graph $G$ and $d>0.$ The induction step follows the setup of the proof of Lemma \ref{K2extend}. Find a clique on the vertices $1,2,\ldots, \omega=\omega(G),$ and define the sets $C_{i,j}$ and $D_i$ as in that proof. For any induced subgraph $H$ of $G$ on $U\subseteq C_{i,j},$ the induced subgraph of $G$ on $U \cup \{1,2,\ldots,j\}$ is isomorphic to $K_{j-2}+(K_2 \cup H).$

If $n-j \geq 2d$ then any induced subgraph $H$ of $C_{i,j}$ on $n-j$ vertices turns into the induced subgraph $K_{j-2}+(K_2 \cup H)$ of $G$ on $n$ vertices. By assumption, the graph $G$ is ${\bf B}_{n,d}$--free, where
$ {\bf B}_{n,d} \supseteq K_{j-2}+(K_2\cup {\bf B}_{n-j,d-1}).$ This shows that $G[C_{i,j}]$ is  ${\bf B}_{n-j,d-1}$--free.

If $n-j < 2d$ then any induced subgraph $H$ of $C_{i,j}$ on $2d$ vertices turns into the induced subgraph $K_{j-2}+(K_2 \cup H)$ of $G$ on $j+2d>n$ vertices. Pick any subgraph $K_{n-2d-2}$ of $K_{j-2}$ to get an induced subgraph $K_{n-2d-2}+(K_2 \cup H)$ of $G.$  By assumption, the graph $G$ is ${\bf B}_{n,d}$--free, where $ {\bf B}_{n,d} \supseteq K_{n-2d-2}+(K_2\cup {\bf B}_{2d,d-1}).$ This shows that $G[C_{i,j}]$ is  ${\bf B}_{2d,d-1}$--free.

Now proceed as in the Proof of Lemma \ref{K2extend} with $\chi(G[D_i])\leq  1 =  g_{\ast,\ast}(1),$ and use Lemma \ref{doubleBumpLemma} to evaluate the sum:
\[\begin{array}{rcl} 
\chi(G) & \leq & \displaystyle  \sum_{i=1} ^\omega \chi(G[D_i]) +  \sum_{1\leq i < j \leq \omega}  \chi(G[C_{i,j}]) \\
& \leq & \displaystyle  \sum_{i=1} ^\omega  g_\ast ( 1) +  \sum_{1\leq i < j \leq \omega}   g_{\max(n-j,2d),d-1}( \omega-j+2 ) \\
& = & \displaystyle  \sum_{1\leq i < j \leq \omega+1}   g_{\max(n-j,2d),d-1}( \omega-j+2 ) \\
& = & \displaystyle  \sum_{j=2}^{\omega+1} (j-1) g_{\max(n-j,2d),d-1}( \omega-j+2 ) \\
& = & \displaystyle  \sum_{k=1}^{\omega} (\omega-k+1) g_{\max(n+k-\omega-2,2d),d-1}( k ) \\
& = & g_{n,d}(\omega).
\end{array}
\]
\end{proof}

\begin{prop}\label{prepp} For $\omega \geq 1$ it holds that
\[ g_{n,d}(\omega) \leq {\omega-1+2d \choose 2d} +  {n-2 \choose 2d+1} \]
and it is sharp exactly when $(d=0, \omega =1, n>2)$ or $(d>0, n>\omega+2d+1).$
\end{prop}
\begin{proof} First we verify the case $d=0,$
\[
{\omega-1+2d \choose 2d} +  {n-2 \choose 2d+1} = {\omega-1 \choose 0} +  {n-2 \choose 1} = 1 + n -2= n-1.
\]
By definition, $g_{n,0}(\omega)=n-1$ for $\omega>1,$ so the proposition statement is an equality in those cases. 
By definition, $g_{n,0}(\omega)=1$ for $\omega=1,$ so the proposition statement is an equality for $n=2$ and a strict inequality for $n>2.$

Now we continue by induction on $d>0.$ Consider that 
\[  \begin{array}{rcl}
g_{n,d}(\omega) & = &  \displaystyle \sum_{k=1}^{\omega} (\omega-k+1) g_{\max(n+k-\omega-2,2d),d-1}( k )  \\
& \leq & \displaystyle \sum_{k=1}^{\omega} (\omega-k+1) \left(    {k-1+2(d-1) \choose 2(d-1)} +  {\max(n+k-\omega-2,2d) -2 \choose 2(d-1)+1}     \right). \\
\end{array} \]
Let's break down the two parts of the sum. The first part is 
\[
\sum_{k=1}^{\omega} (\omega-k+1)   {k-1+2(d-1) \choose 2(d-1)}  = \sum_{i=0}^{\omega-1} ((\omega-1)+1-i)   {i+2d-2 \choose 2d-2}  = { \omega - 1 +2d \choose 2d }
\]
by Lemma \ref{doubleBumpLemma}. The second part is 
\[
\sum_{k=1}^{\omega} (\omega-k+1)  {\max(n+k-\omega-2,2d) -2 \choose 2(d-1)+1}.
\]
The binomial coefficient vanishes unless $n+k-\omega-2 \geq 2d+1,$ or equivalently, that $k \geq \omega - n + 2d + 3.$ The second part may be rewritten as:
\[
\begin{array}{rcl}
\displaystyle  \sum_{k=\omega - n + 2d + 3}^{\omega} (\omega-k+1)  {n+k-\omega-4 \choose 2d-1} &  = & \displaystyle  \sum_{i=0}^{n - 2d - 3} ((n - 2d - 3)+1-i)  {i + 2d -1 \choose 2d-1}    \\
    & = & \displaystyle  {n - 2d - 3 + 2d -1 +2  \choose 2d-1 + 2 }    \\
    & = & \displaystyle  {n -2   \choose 2d+1 },    \\
\end{array}
\]
with an application of Lemma \ref{doubleBumpLemma}. Adding the two pieces up proves inequality.

Now to the question of when the inequality is an equality or sharp. For $d=0$ it's sharp exactly when $\omega=1$ and $n>2.$ For $d=1,$ in the expansion of $g_{n,1}(\omega),$ the terms are of the type:
\[ g_{\max(n+k-\omega-2,2d),d-1}( k ) \]
for $1 \leq k \leq \omega .$ The only one that could be sharp is for $k=1:$
\[ g_{\max(n-\omega-1,2),0}( 1 ) \]
The condition is that $n-\omega-1>2,$ or equivalently,  $n>\omega+3 $ for it to be sharp. 

For $d>1$ we prove it by induction on $d.$ The inequality for $g_{n,d}(\omega)$ is sharp exactly when $n> \omega +2d+1.$ We expanded $g_{n,d}(\omega)$ into terms of the type
\[ g_{\max(n+k-\omega-2,2d),d-1}( k ). \]
By induction, the inequality for one of them is sharp exactly when
\[
\max(n+k-\omega-2,2d) > k + 2(d-1) + 1.
\]
It can never be that $2d> k + 2(d-1) + 1$, as $k\geq 1.$ It remains that
\[ n+k-\omega-2>2d \,\, \textrm{ and } \,\, n+k-\omega-2> k + 2(d-1) + 1. \]
The rightmost inequality implies the leftmost one, as $k\geq 1.$ It simplifies to
the desired
\[ n> \omega +2d+1
\]
and we are done.
\end{proof}


 
 \begin{prop}\label{gTriangleFree} $ g_{n,d}(2) = n-1$ for $d\geq 0$ and $n\geq 2(d+1).$
\end{prop}
\begin{proof} We prove it by induction on $d.$ For $d=0$ it is true by definition. Assume that $d>0.$
By the assumption $n\geq 2(d+1),$
\[ \max(n-2,2d) = n-2. \] 
By definition,
\[ g_{n,d}(2)  = 2g_{\max(n-3,2d),d-1}(1) + g_{\max(n-2,2d),d-1}(2) = 1 + g_{n-2,d-1}=1+n-2=n-1.\]
\end{proof}



 We record the exact values of $g_{n,d}(\omega)$ for some small $n,d.$ The proposition above applies exactly for $d>0$ when $\omega \geq n - 2d - 1 .$
 
\begin{exa}  For $\omega > 1$ the functions 
 \[\begin{array}{|c|c|c|c|c|c|}
 \hline
 g_{2,0}(\omega) & g_{3,0}(\omega) & g_{4,0}(\omega) & g_{5,0}(\omega) & g_{6,0}(\omega) & g_{7,0}(\omega) \\
 \hline & g_{4,1}(\omega)& g_{5,1}(\omega)  & g_{6,1}(\omega)  & g_{7,1}(\omega) & g_{8,1}(\omega)  \\
  \hline && g_{6,2}(\omega)& g_{7,2}(\omega)  & g_{8,2}(\omega)  & g_{9,2}(\omega) \\
   \hline &&& g_{8,3}(\omega)  & g_{9,3}(\omega)  & g_{10,3}(\omega) \\
    \hline &&&& g_{10,4}(\omega)  & g_{11,4}(\omega) \\
     \hline &&&&& g_{12,5}(\omega) \\
 \hline
 \end{array} \]
 are
 \[\begin{array}{|c|c|c|c|c|c|}
 \hline
 \omega \mapsto 1 &  \omega \mapsto 2 &  \omega \mapsto 3& \omega \mapsto 4 &  \omega \mapsto 5 &  \omega \mapsto 6  \\
 \hline &  \omega \mapsto {\omega+1 \choose 2}  &  \omega \mapsto {\omega+1 \choose 2}+1 &  2 \mapsto 5  & 2 \mapsto 6 & 2 \mapsto 7 \\
 &   &  &  \omega \mapsto {\omega+1 \choose 2}+4  & 3 \mapsto 13& 3 \mapsto 16  \\
 &   &  &  \textrm{for $\omega \geq 3$}  &  \omega \mapsto {\omega+1 \choose 2}+10 & 4 \mapsto 26  \\
  &   &  & &   \textrm{for $\omega \geq 4$}  &  \omega \mapsto {\omega+1 \choose 2}+20  \\
   &   &  & & &   \textrm{for $\omega \geq 5 $}    \\
  \hline && \omega \mapsto {\omega+3 \choose 4} & \omega \mapsto {\omega+3 \choose 4}+1  & 2 \mapsto 7 & 2 \mapsto 8 \\
    &&&  & \omega \mapsto {\omega+3 \choose 4}+6  & 3 \mapsto 26  \\
     &&&  & \textrm{for $\omega \geq 3 $} &  \omega \mapsto {\omega+3 \choose 4}+21    \\
      &&&&  & \textrm{for $\omega \geq 4 $}   \\
   \hline &&&  \omega \mapsto {\omega+5 \choose 6}  & \omega \mapsto {\omega+5 \choose 6} +1   & 2\rightarrow 9   \\
   &&&   &  &  \omega \mapsto {\omega+5 \choose 6} + 8 \\
    &&&&  & \textrm{for $\omega \geq 3 $}   \\
    \hline &&&& \omega \mapsto {\omega+7 \choose 8}    & \omega \mapsto {\omega+7 \choose 8} +1 \\
     \hline &&&&&  \omega \mapsto {\omega+9 \choose 10}  \\
 \hline
 \end{array} \]
 where the only values not given by the results above are
 \[ g_{7,1}(3) =  3 g_{3,0}(1) + 2 g_{4,0}(2) +  g_{5,0}(3)  =  3\cdot 1 + 2 \cdot 3 + 4 = 3+6+4 = 13, \]
 \[ g_{8,1}(3) =  3 g_{4,0}(1) + 2 g_{5,0}(2) +  g_{6,0}(3)  = 3\cdot 1 + 2 \cdot 4 + 5 =  16, \]
 \[ g_{8,1}(4) =  4 g_{5,0}(1)  + 3 g_{4,0}(2) + 2 g_{5,0}(3) +  g_{6,0}(4) = 4\cdot 1 + 3 \cdot 3 + 2\cdot 4 + 5 = 26, \]
 \[ g_{9,2}(3) = 3 g_{5,1}(1) + 2 g_{6,1}(2) + g_{7,1}(3)= 3\cdot 1 +2\cdot 5 + 13 = 26. \]
 \end{exa}
 
\section{Algebraic results}

The \emph{independence complex} $\mathrm{Ind}(G)$ of a graph $G$ is the simplicial complex of independent sets of $G.$ The following three lemmas are elementary and essentially well-known:

\begin{lem} \label{indH0} $\dim \tilde{H}_0(\mathrm{Ind}(G)) = 0$ if and only if the complement of $G$ is connected.
\end{lem}

\begin{lem} \label{indSusp} $\dim \tilde{H}_{d+1}(\mathrm{Ind}(K_2 \cup G)) = \dim \tilde{H}_d(\mathrm{Ind}(G))$ for $d\geq 0.$
\end{lem}

\begin{lem} \label{indDisj} $\dim \tilde{H}_{d}(\mathrm{Ind}(K_m + G)) = \dim \tilde{H}_d(\mathrm{Ind}(G))$ for $d > 0.$
\end{lem}

This is central for this paper:

\begin{lem}\label{central1}
If $H \in {\bf B}_{n,d}$ then $\dim \tilde{H}_d(\mathrm{Ind}(H)) > 0.$
\end{lem}

\begin{proof}
First the case $d=0.$ The graphs in ${\bf B}_{n,0}$
are defined by that their maximal degree is $n-1.$ Let $v$ be a vertex of $H$ of degree $n-1.$ The complement of $H$ is disconnected as $v$ is isolated in it. By Lemma \ref{indH0}, $\dim \tilde{H}_0(\mathrm{Ind}(H)) > 0.$

Now the case $d>0$ by induction. By definition
\[ H \in {\bf B}_{n,d} = \bigcup_{m=2d}^{n-2} K_{n-m-2}+(K_2\cup {\bf B}_{m,d-1}). \]
Pick the $m$ with $H \in K_{n-m-2}+(K_2\cup {\bf B}_{m,d-1}).$ Then $H = K_{n-m-2}+(K_2\cup H')$ for some $H'\in {\bf B}_{m,d-1}.$ By induction, $\dim \tilde{H}_{d-1}(\mathrm{Ind}(H')) > 0.$ By Lemma \ref{indSusp}, 
$\dim \tilde{H}_{d}(\mathrm{Ind}(K_2\cup H')) > 0.$ Finally, by Lemma \ref{indDisj},
\[
\dim \tilde{H}_d(\mathrm{Ind}(H)) = \dim \tilde{H}_{d}(\mathrm{Ind}(K_{n-m-2}+(K_2\cup H')))>0.
\]
\end{proof}

\begin{lem} \label{central2}
If $G$ is a graph with $\dim \tilde{H}_{d}(\mathrm{Ind}(H)) = 0$ for every induced subgraph $H$ on $n$ vertices of $G,$ then $G$ is a ${\bf B}_{n,d}$--free graph.
\end{lem}
\begin{proof}
Assume to the contrary that $G$ is not a ${\bf B}_{n,d}$--free graph, but\linebreak $\dim \tilde{H}_d(\mathrm{Ind}(H)) = 0$ for every induced subgraph $H$ on $n$ vertices of $G.$
Let $H$ be an induced subgraph of  $G$  on $n$ vertices that is not ${\bf B}_{n,d}$--free. All graphs in ${\bf B}_{n,d}$ are on $n$ vertices, so $H \in {\bf B}_{n,d}.$
By Lemma \ref{central1},  $\dim \tilde{H}_d(\mathrm{Ind}(H)) > 0$ and we have a contradiction.
\end{proof}

\begin{defn}
\[
\beta_{i,j}(I_G) = \sum_{W \in { V(G) \choose j } } \dim_{\mathbf k}\tilde{H}_{j-i-2}(\mathrm{Ind}(G[W]))
\]
\end{defn}

\begin{thm} \label{mainThm}
The class of graphs with $ \beta_{i,j}(I_G) = 0$ is $\chi$--bounded by $g_{j,j-i-2}(\omega).$
\end{thm}
\begin{proof}
If $ \beta_{i,j}(I_G) = 0 $ then 
$ \dim_{\mathbf k}\tilde{H}_{j-i-2}(\mathrm{Ind}(G[W])) = 0 $ for all $W \in { V(G) \choose j }.$
So, by Lemma \ref{central2}, the graph $G$ is a ${\bf B}_{j,j-i-2}$--free graph.
By Theorem \ref{bigthm}, the class of  ${\bf B}_{j,j-i-2}$--free graphs is $\chi$--bounded by $g_{j,j-i-2}(\omega).$
\end{proof}

\begin{cor} \label{mainCor1} The class of graphs with $ \beta_{i,j}(I_G) = 0$ is $\chi$--bounded by
\[ {\omega-1+2(j-i-2) \choose 2(j-i-2)} +  {j-2 \choose 2(j-i-2)+1}. \] 
\end{cor}
\begin{proof} Use Theorem \ref{mainThm} and  Proposition \ref{prepp}. 
\end{proof}

\begin{cor} \label{mainCor2} If $G$ is a triangle-free graph with $\beta_{i,j}(I_G) = 0$ then $\chi(G) \leq j-1.$
\end{cor}
\begin{proof} Use Theorem \ref{mainThm} and  Proposition \ref{gTriangleFree}. 
\end{proof}



\section{Asymptotically $\chi$-bounded graphs}

Reed and Yuditsky \cite{RY25} defined an asymptotic version:
\begin{defn}
An hereditary class $\mathcal{C}$ of graphs is \emph{asymptotically $\chi$-bounded by~$f$} if
\[
\lim_{n \rightarrow \infty} {\mathbf P}  [ \chi(G)\leq f(\omega(G))  \mid \textrm{$G$ is an order $n$ graph in the class $\mathcal{C}$}] = 1.
\]
\end{defn}
Towards the Gyárfás-Sumner conjecture, they proved that for any tree $T,$ the class of $T$--free graphs is asymptotically $\chi$-bounded by $\omega.$

\begin{defn}
A betti number $\beta_{i,j}(I_G)$ is \emph{parabolic} if $(j-i)^2 \geq j+i+2.$
\end{defn}

This  theorem was proved by Engström and Orlich \cite{EO23}.

\begin{thm}\label{eo} Let $\beta_{i,j}(I_G)$ be a parabolic betti number with $j\geq i+3.$ Then, for almost every graph with $ \beta_{i,j}(I_G) = 0$ we have:
\begin{itemize}
\item[1.] $\beta_{k,l}(I_G)=0$ whenever $l-k=j-i,$ and
\item[2.] parabolic $\beta_{k,l}(I_G)>0$ whenever  $l-k<j-i.$
\end{itemize}
\end{thm}

\begin{cor} \label{asymBound} If  $(j-i)^2 \geq j+i+2$ and $j-i\geq 3,$ then the hereditary class of graphs with $ \beta_{i,j}(I_G) = 0$ is 
asymptotically $\chi$-bounded by 
\[ {\omega-1+2(j-i-2) \choose 2(j-i-2)}. \]
\end{cor}
\begin{proof}
Consider $\beta_{j-i-2,2j-2i-2}(I_G).$ We have that $(2j-2i-2)-(j-i-2)=j-i,$ so by Theorem \ref{eo} above $\beta_{j-i-2,2j-2i-2}(I_G)=0$ for almost all $G.$
By Corollary  \ref{mainCor1},
\[
\begin{array}{rcl}
\chi(G) & \leq &  \displaystyle {\omega-1+2((2j-2i-2)-(j-i-2)-2) \choose 2((2j-2i-2)-(j-i-2)-2)} +  {2j-2i-2-2 \choose 2((2j-2i-2)-(j-i-2)-2)+1} \\
& =  & \displaystyle  {\omega-1+2(j-i-2) \choose 2(j-i-2)} +  {2j-2i-2-2 \choose 2(j-i-2)+1} \\
& =  & \displaystyle  {\omega-1+2(j-i-2) \choose 2(j-i-2)} +  {2j-2i-4 \choose 2j-2i-3} \\
& =  & \displaystyle  {\omega-1+2(j-i-2) \choose 2(j-i-2)}. \\
\end{array}
\]
\end{proof}
\section{Complexity theoretic results}

In this section the time complexity for algorithms and subroutines is many times essentially $O(m)$ for graphs with $n$ vertices and $m$ edges. But to be more precise, it is actually $O(m+n)$ unless
certain corner cases are avoided. Unfortunately those corner cases are common in subroutines. Furthermore, working with $O()$--notation turns out to be unsuitable in recursions due to that statements like $O(1)+O(2)+\cdots+O(t)$ does not make sense. Usually we will prove that algorithms run faster than $C\varphi(G)(m+n)$ where $C$ is a positive constant, $G$ is a non-empty graph and $\varphi(G)$ is a positive graph invariant.

All graph coloring in this text so far builds on repeated applications of Lemma~\ref{K2extend}. The main complexity problem is that it requires one to find a maximum clique, and that is notoriously slow.
The first goal of this section is to replace the maximum clique in the proof of Lemma~\ref{K2extend} by something weaker that is faster to find. A clique on $K=\{1,2,\ldots, k\}$ in $G$ is selected with 
$k=\omega(G)$ to achieve the following two facts in the proof:
\begin{itemize}
\item[(1)] There is no $v \in V(G)\setminus K$ adjacent to all vertices of $K.$
\item[(2)] There is no $u\in K$ together with two adjacent vertices $v ,w \in V(G)\setminus K$ that are adjacent to all vertices of $K\setminus u.$
\end{itemize}
A clique only satisfying condition (1) is maximal.

Andrade, Resende and Werneck  \cite{arw12}  found a fast heuristic to find large independent sets by local improvements. As mentioned in their paper, it is also valid to find cliques, which is the relevant context for this paper. Transforming a clique by deleting $k-1$ vertices and then adding $k$ vertices to get a new larger clique is called a \emph{$k$--improvement}. If no $k$--improvement is possible, the clique is \emph{$k$--maximal}. To be 1-maximal is equivalent to ordinary maximal. It is easy to prove that:
\begin{lem}\label{1-imp}
There is a positive constant $C_{\mathrm{1-imp}}$ such that given any clique in a non-empty graph on $n$ vertices and $m$ edges one can find an 1-improvement (or prove that none exist) in $C_{\mathrm{1-imp}}(m+n)$ time.
\end{lem}
The heuristic by Andrade et al is fast in practice and they also proved that:
\begin{lem}\label{2-imp}
There is a positive constant $C_{\mathrm{2-imp}}$ such that given any 1-maximal clique in a non-empty graph on $n$ vertices and $m$ edges one can find an 2-improvement (or prove that none exist) in $C_{\mathrm{2-imp}}(m+n)$ time.
\end{lem}
Putting them together we may find a 2-maximal clique.
\begin{lem}\label{comp1}
There is a positive constant $C_{\mathrm{2-max}}$ such that given any non-empty graph on $n$ vertices and $m$ edges with maximum clique size $\omega$ one may find a 2-maximal clique in 
$C_{\mathrm{2-max}}\omega(m+n)$ time.
\end{lem}
\begin{proof}
Start with an empty clique and use Lemma~\ref{1-imp} and  Lemma~\ref{2-imp} until the clique is 2-maximal. The number of applications of them is in total at most $\omega,$ as that is the maximum size of a clique in $G.$ Set $C_{\mathrm{2-max}} = \max(C_{\mathrm{1-imp}} ,C_{\mathrm{2-imp}}).$
\end{proof}
\begin{lem}
If 
\begin{itemize}
\item[-] ${\bf H}$ is a set of graphs,
\item[-]  the class of ${\bf H}$--free graphs is  $\chi$-bounded by $f_{\bf H},$ and
\item[-]  there exist a positive constant $C_{\bf H}$ such that given any non-empty graph~$G$ on $n$ vertices and $m$ edges with maximum clique size $\omega$ one may find a
$f_{\bf H}(\omega)$--coloring of $G$ in $C_{\bf H}\omega(m+n)$ time,
\end{itemize}
then
\begin{itemize}
\item[-] the class of  $K_2 \cup {\bf H  } $--free graphs is $\chi$-bounded by
\[
f_{K_2 \cup {\bf H  } }(\omega) = \sum_{k=1}^{\omega}  (\omega-k+1)  f_{{\bf H  } }(k),
\]
\item[-] and there exist a positive constant $C_{K_2 \cup {\bf H  } }$ such that given any non-empty graph~$G$ on $n$ vertices and $m$ edges with maximum clique size $\omega$ one may find a
$f_{K_2 \cup {\bf H  } }(\omega)$--coloring of $G$ in $C_{K_2 \cup {\bf H  } }\omega(m+n)$ time.
\end{itemize}
\end{lem}
\begin{proof}
The statements on $\chi$--boundedness are all in Lemma~\ref{K2extend}. The only new thing in Lemma~\ref{comp1} are the complexity theoretic statements. We follow the proof of Lemma~\ref{K2extend} and work out the time used. This starts off with finding a 2-maximal clique instead of a maximum clique. It is found in time $C_{\mathrm{2-max}}\omega(m+n)$ according to Lemma~\ref{comp1}.

Let $\tilde\omega$ be the cardinality of the 2-maximal clique. The next step is to find the sets $C_{i,j}$ for $1\leq i < j \leq \tilde\omega$ and $D_i$ for $1\leq i \leq \tilde\omega.$ It is straightforward to do that in $C'\omega(m+n)$ time, for some constant $C'.$ Actually a minor adaptation of the algorithm to find the 2-maximal clique provides this information almost for free.

Each $D_i$ is independent, and coloring all of them is linear in the number of vertices, say $C''n.$

Now let $n_{i,j}$ and $m_{i,j}$ be the number of vertices and edges in each $C_{i,j},$ respectively. Coloring all of them is done in
\[
\sum_{1\leq i < j \leq \tilde\omega}  C_{\bf H}\omega(m_{i,j}+n_{i,j}) \leq C_{\bf H}\omega(m+n)
\]
time. In total at most $C_{K_2 \cup {\bf H  } }\omega(m+n)$ time is used, where essentially one may set $C_{K_2 \cup {\bf H  } } = C_{\mathrm{2-max}} + C' + C'' + C_{\bf H}.$
\end{proof}

A minor detail above is that if  $\tilde\omega<\omega$ then a better chromatic bound is actually achieved:
\[ \sum_{k=1}^{\tilde \omega}  (\omega-k+1)  f_{{\bf H  } }(k) \]

In the same way as Lemma~\ref{comp1} is the complexity version of Lemma~\ref{K2extend}, the following Theorem~\ref{comp2} is the complexity version of Theorem~\ref{bigthm}.
\begin{thm}\label{comp2}
For any integers $d\geq 0$ and $n\geq 2(d+1)$ there is a constant $C$ such that every ${\bf B}_{n,d}$--free graph $G$ on $n'$ vertices and $m'$ edges is 
$g_{n,d}(\omega(G))$--colorable in $C\omega(G)(m'+n')$ time.
\end{thm}
\begin{proof}
The existence of the coloring is the content of Theorem~\ref{bigthm}. To be proved is the complexity result. In the proof of Theorem~\ref{bigthm} the initial part is that graphs of constant maximal degree are colored optimally. That can be done in $C'(m'+n')$ time using a first-fit coloring. After that, a version of Lemma~\ref{K2extend} with refined bookkeeping is applied $d$ times. From Lemma~\ref{comp1} we know that each application is done faster than $C''\omega(G)(m'+n')$ for some $C''.$ Set $C=C'+dC''.$
\end{proof}

We also get complexity versions of Corollaries \ref{mainCor1} and \ref{mainCor2}.

\begin{cor}\label{n3} Let $i$ and $j$ be positive integers. Any graph $G$ on $n$ vertices and clique number $\omega$ with $\beta_{i,j}(I_G) = 0$ has chromatic number
\[ \chi \leq {\omega-1+2(j-i-2) \choose 2(j-i-2)} +  {j-2 \choose 2(j-i-2)+1} \]
and such a coloring may be found in $O(n^3)$ time. 
\end{cor}
\begin{proof}
The coloring result in Corollary \ref{mainCor1} follows from Theorem~\ref{bigthm}, whose complexity version is Theorem~\ref{comp2} above. According to that, the coloring is found in $C\omega(m+n)$ time where $m$ is the number of edges in $G$ and $\omega$ is the clique number. From $m\leq n^2$ and $\omega\leq n$ it follows that $O(n^3)$ time is used.
\end{proof}
In the same manner:
\begin{cor}\label{n2} Let $i$ and $j$ be positive integers. Any triangle-free graph $G$ on $n$ vertices with $\beta_{i,j}(I_G) = 0$ 
is $(j-1)$--colorable and such a coloring may be found in $O(n^2)$ time.
\end{cor}


\section{Further considerations}
\subsection{Other classes of forbidden graphs}
Instead of the ${\bf B}_{n,d}$ sets defined in Section 3, one may consider other sets of forbidden graphs as they do not contain all possible obstructions to vanishing syzygies. For example, $C_5$ is a possible addition to ${\bf B}_{5,1}.$ After adding $C_5$ to ${\bf B}_{5,1},$ the recursion
\[{\bf B}_{n,d} = \bigcup_{m=2d}^{n-2} K_{n-m-2}+(K_2\cup {\bf B}_{m,d-1}) \] 
is still valid and may be employed to extend ${\bf B}_{n,d}$ for $d\geq 2$ and $n\geq 5+2(d-1).$
Moreover, the complement of any cycle $C_n$ may be added to  ${\bf B}_{n,1}$ for $n>5.$

Another aspect would be to reduce the size of ${\bf B}_{n,d}$ and sacrifice the bound achieved to get an easier proof. There are at least two natural candidates in that direction. The first one is to set ${\bf B}_{n,0}' = {\bf B}_{n,0}$ and
\[{\bf B}_{n,d}' = K_2\cup {\bf B}_{n-2,d-1}'. \]
The ${\bf B}_{n,d}'$--free graphs are $\chi$--bounded by a polynomial of the same degree as the ${\bf B}_{n,d}$--free graphs, but with much worse coefficients. On the other hand, Theorem~\ref{bigthm} may be replaced with the easier Corollary~\ref{corBoundingSimple}. Completely disregarding the coefficients and only aiming to preserve the degree of the $\chi$--bounding polynomials while simplifying the argument even more, one may replace Corollary~\ref{corBoundingSimple} by the basic inductive step in Wagon's proof of that $pK_2$--free graphs are $\chi$--bounded by a polynomial of degree $2p-2.$

A further reduction would be to start off with ${\bf B}_{n,0}'' = \{ \textrm{the star on $n$ vertices} \}$ and continue with ${\bf B}_{n,d}'' = K_2\cup {\bf B}_{n-2,d-1}''.$ The upside is that each set only contains one graph:
\[ {\bf B}_{n,d}'' =  \{ (d+1)K_2 \cup \textrm{the star on $n-2(d+1)$ vertices} \}. \]
The main downside of this construction is that a star on $m$ vertices is $\chi$--bounded by $k\frac{\omega^{m-2}}{\log^{m-3} (\omega)}$ for some constant $k.$ %See for example the introduction of Polynomial χ-binding functions for t-broom-free graphs Xiaonan Liu ∗ Joshua Schroeder † Zhiyu Wang ‡ Xingxing Yu § November 30, 2022
This shows that ${\bf B}_{n,0}''$--free graphs are $\chi$--bounded by $k\frac{\omega^{n-2}}{\log^{n-3} (\omega)},$ while the ${\bf B}_{n,0}'$--free graphs are $\chi$--bounded by $n.$ This factor propagates to that
the ${\bf B}_{n,d}''$--free graphs are $\chi$--bounded by $\frac{k}{n-2(d+1)}\frac{\omega^{n-2(d+1)-2}}{\log^{n-2(d+1)-3} (\omega)}$ times whatever the ${\bf B}_{n,d}'$--free graphs are $\chi$--bounded by. Thus, the ${\bf B}_{n,d}''$--free graphs also proves polynomially $\chi$--bounded, but by a worse degree.

A disjoint union of stars is called a \emph{star-forest}. It was proved by Scott, Seymour and Spirkl \cite{sss} that for any star-forest $H,$ the $H$--free graphs are polynomially $\chi$--bounded. They did not make any effort in optimizing degrees of the polynomials, and it does not seem straight-forward to use their result right off to improve on our results. What perhaps would be feasible is to exclude many star-forests at once. Define
\[
{\bf B}_{n,d}^\dagger = \{ \textrm{star-forests on $d+1$ components and $n$ vertices}  \}.
\]
Clearly the ${\bf B}_{n,d}^\dagger$--free graphs are polynomially $\chi$--bounded as every forbidden graph is a star-forest. It's an interesting research problem to find a better chromatic bound for the ${\bf B}_{n,d}^\dagger$--free graphs. The independence complex of a star-forest on $d+1$ components collapses onto the $d$--dimensional boundary of a cross polytope, providing the syzygy obstruction.

\subsubsection*{Acknowledgements}
The author would like to thank Ignacio García-Marco, Philippe Gimenez and Kolja Knauer.

\begin{thebibliography}{100} 

\bibitem{arw12}
Diego V. Andrade, Mauricio G. Resende and Renato W.  Werneck. Fast local search for the maximum independent set problem. \emph{Journal of Heuristics} 18, 525--547 (2012)

\bibitem{BC18}
Arpitha P. Bharathi and Sheshayya A. Choudum. Colouring of $(P_3 \cup P_2)$-free graphs. \emph{Graphs and Combinatorics} 34, 97--107 (2018)

\bibitem{BT25}
Romain Bourneuf and St\'ephan Thomass\'e. Bounded twin-width graphs are polynomially $\chi$-bounded. \emph{Advances in Combinatorics}. (2025)

\bibitem{REF1}
Marcin Briański, James Davies and Bartosz Walczak. Separating polynomial $\chi$-boundedness from $\chi$-boundedness. \emph{Combinatorica} 44, 1--8 (2024)

\bibitem{CS23}
 Maria Chudnovsky and Paul Seymour. Even-hole-free graphs still have bisimplicial vertices. \emph{Journal of Combinatorial Theory B} 161, 331--381 (2023)

\bibitem{CS19}
Maria Chudnovsky and Vaidy Sivaraman. Perfect divisibility and 2-divisibility. \emph{Journal of Graph Theory} 90, 54--60 (2019)

\bibitem{Dav22} 
James Davies. Improved bounds for colouring circle graphs. \emph{Proceedings of the American Mathematical Society} 150, 5121–5135 (2022)
 
\bibitem{DM21}
James Davies and Rose McCarty. Circle graphs are quadratically $\chi$-bounded. \emph{Bulletin of the London Mathematical Society} 53, 673--679 (2021)
 
\bibitem{de}
Anton Dochtermann and Alexander Engström. Algebraic Properties of Edge Ideals via Combinatorial Topology. \emph{The Electronic Journal of Combinatorics} 16 (2009)

\bibitem{EO23}
Alexander Engström and Milo Orlich. The regularity of almost all edge ideals. \emph{Advances in Mathematics} 435 (2023)

\bibitem{fro}
Ralf Fr\"oberg, On Stanley–Reiner rings, Topics in Algebra, vol. 2, Banach Center Publ., Warsaw (1990)

\bibitem{ran}
Hugues Randriambololona. The syzygy distinguisher. \emph{Eurocrypt 2025} 324--354 (2025)

\bibitem{RY25}
Bruce Reed and Yelena Yuditsky. The asymptotic $\chi$-boundedness of hereditary families. arXiv:2506.01070 (2025)

\bibitem{sss}
Alex Scott, Paul Seymour and Sophie Spirkl. Polynomial bounds for chromatic number II. Excluding a star-forest. \emph{Journal of Graph Theory}  101, 2, 318--322 (2022)

\bibitem{S74}
Dieter Seinsche. On a property of the class of $n$-colorable graphs. \emph{Journal of Combinatorial Theory B} 16, 191--193 (1974)

\bibitem{W80}
Stanley Wagon. A bound on the chromatic number of graphs without certain induced subgraphs. \emph{Journal of Combinatorial Theory B} 29, 345--346 (1980)

\end{thebibliography}
\end{document}

 
